\hspace{3mm}

\centerline{\bf \Huge Домашнее задание.}
\centerline{\bf \Huge Графы}

\begin{flushright}
    \scriptsize{Авраам Линкольн дал людям свободу,\\ а полковник Кольт уравнял их шансы}
\end{flushright}

\problem{1!}Могут ли степени вершин в графе быть равны:

а) $100, 99, 98, \dotsc, 3, 2, 1$?
  
б) $9, 9, 8, 8, 6, 6, 4, 4, 1, 1$
  
в) $6, 6, 6, 5, 5, 3, 2, 2?$

\problem{2!} В графе степени всех вершин равны 3 и между любыми двумя вершинами существует путь длиной не более 2. Какое наибольшее число вершин может быть в этом графе?

\problem{3} В стране Калмыстан есть несколько городов, некоторые из них соединены между собой авиалиниями, причем известно, что из любого города можно попасть в любой другой(возможно с пересадками) и нет зацикливающихся путей по авиалиниям. Докажите, что некоторые города можно объявить Республикой Калмыкия, а остальные Республикой Дагестан так, что все перелеты в стране будут межреспубликанскими.

\problem{4}В Дагестане 2024 мегополиса, некоторые соединены между собой сверхзвуковой дорогой, причём из каждого мегаполиса можно добраться до любого другого. Докажите, что можно несколько дорог из сверхзвуковых сделать сверхсветовыми так, чтобы из каждого города выходило нечётное число сверхсветовых дорог.